\begin{document}
\chapter*{د دې نسخې يادښتونه}
دا يادښتونه د ژباړې د نسخې له خوا زيات شوي دي؛ د اصلي متن برخه نۀ
دي۔ په اصلي متن کښې نښې پر خپل ځاے ساتل شوې دي۔

په \olref[sfr][rel][set]{sec} کښې د طبيعي عددونو د ترتيبي اړيکو په
بېلګه کښې $I$ بې له ښکاره تعريفه کارول شوے دے۔ له سياقه داسې
ښکاري چې مطلب ئې د طبيعي عددونو د عينيت اړيکه ده۔

په \olref[sfr][rel][tre]{sec} کښې د څانګې د تعريف په فورمول کښې
$X$ راغلے، خو د ونې د غړو سټ هلته $A$ نومول شوے دے۔ له سياقه
داسې ښکاري چې په دې فورمول کښې هم مطلب $A$ دے۔ دا د اصلي متن د
نښې ناڅرګندتيا ده۔

په \olref[sfr][rel][ord]{sec} او \olref[sfr][rel][ops]{sec} کښې
$R^+$ د بېلو جوړښتونو دپاره کارول کېږي: په لومړۍ برخه کښې د قطر
په ورزياتولو انعکاسي بندښت ښيي، او په دويمه کښې انتقالي بندښت۔
په هره برخه کښې بايد د هماغې برخې تعريف تعقيب شي۔

په \olref[sfr][fun][bas]{examplefunext} کښې اصلي متن ورکړل شوے
عدد $n$ نوموي، خو بيا د هغه پر ځاے $x$ کاروي۔ ژباړه د هماغه
ننوت دپاره نوم نه تکراروي او فورمول پر خپل ځاے ساتي
(OLFUN-003)۔ په همدې برخه کښې انګرېزي سرچينه د مربع جذر په
بېلګه کښې «مثبت» جذر وايي؛ د صفر دپاره جذر صفر دے، نو ژباړه
د ټولو طبيعي عددونو دپاره «غير منفي (اصلي)» جذر وايي
(OLFUN-002)۔ د مثبتو صحيح عددونو د دوو جذرونو خبره نۀ ده بدله شوې۔

په \olref[sfr][fun][inv]{sec} کښې د کيڼ معکوس د شتون په ثبوت کښې
يو $a \in A$ ټاکل کېږي۔ دا ګام د $A$ ناتش والے غواړي۔ که
$A=\varnothing$ او $B\ne\varnothing$ وي، تشه تابع له $A$ څخه $B$
ته يو پر يو ده، خو له $B$ څخه تش سټ ته تابع نشته، نو کيڼ معکوس
نۀ لري۔ که دواړه سټونه تش وي، تشه تابع خپل معکوس ده۔ د اصلي
متن قضيه دا استثنا په ښکاره نۀ بيانوي۔ په ژباړل شوې قضيه کښې
د ناتشې تعريف ساحې ساده شرط ورزيات شوے دے۔ په بشپړ عموميت کښې
يوه يو پر يو تابع هله کيڼ معکوس لري چې $A\ne\varnothing$ وي يا
$B=\varnothing$ وي (OLFUN-001)۔

په \olref[sfr][fun][rel]{sec} کښې انګرېزي سرچينه د تابعې ګراف
«پر $A\times B$ اړيکه» بولي۔ د همدې کتاب د تعريف له مخې، کره خبره
دا ده چې دا د $A$ او $B$ ترمنځ اړيکه، يعنې د $A\times B$ فرعي
سټ، دے۔ ژباړه همدا کره عبارت کاروي (OLFUN-004)۔

په \olref[sfr][rel][ops]{sec} کښې تر $C$ پورې د اړيکې محدودونه
دواړه مختصات په $C$ کښې غواړي؛ خو په \olref[sfr][fun][rel]{sec}
کښې تر $C$ پورې د تابعې محدودونه يوازې ننوت په $C$ کښې غواړي۔
نو دا دوه تعريفونه په عمومي ډول يو شان نۀ دي۔ د بېلګې په توګه،
که $f(0)=1$ او $C=\{0\}$ وي، نو د تابعې محدودونه جوړه
$\tuple{0,1}$ ساتي؛ د اړيکې هغه محدودونه ئې نۀ ساتي۔ هر تعريف
بايد د خپلې برخې له مخې ولوستل شي۔ د تابعې ښکاره تعريف سم دے؛
يوازې له مخکيني تعريف سره د عينيت ادعا محدوده شوې ده (OLFUN-005)۔
\end{document}
